\documentclass[11pt,a4paper]{article}

\usepackage[utf8]{inputenc}
\usepackage[T1]{fontenc}
\usepackage[ngerman]{babel}
\usepackage[a4paper,left=22mm,right=22mm,top=18mm,bottom=20mm]{geometry}
\usepackage{xcolor}
\usepackage{hyperref}

% LBBW-inspirierte Farben und Schrift passend zum Lebenslauf.
\definecolor{CVPrimary}{HTML}{123250}
\definecolor{CVAccent}{HTML}{37c391}

\hypersetup{colorlinks=true,urlcolor=CVPrimary}
\renewcommand{\familydefault}{\sfdefault}
\pagestyle{empty}
\setlength{\parindent}{0pt}
\setlength{\parskip}{0.65em}
\setlength{\emergencystretch}{2em}

% Für ein festes Versanddatum beispielsweise: 14. Oktober 2026
\newcommand{\Briefdatum}{\today}

\begin{document}
\color{black}

{\raggedleft \Briefdatum\par}
\vspace{5mm}
{\large\bfseries\textcolor{CVPrimary}{Bewerbung als Trainee Informationstechnologie (m/w/d)}}\par
{\small CIO Office \& Project Management | Referenznummer 17859 | Start 01.04.2027}\par
{\color{CVAccent}\rule{\linewidth}{0.6pt}}\par
\vspace{1mm}

Sehr geehrte Damen und Herren,

Nur durch Innovation bleibt eine Bank wettbewerbsfähig, und Innovation entsteht heute vor allem in der IT. Genau hier möchte ich im IT-Traineeprogramm der LBBW mitwirken. Besonders reizt mich, dass ich verschiedene Themenfelder der IT kennenlerne und von Beginn an einen Beitrag zur IT-Delivery leisten kann.

Im dualen Bachelorstudium der Wirtschaftsinformatik bei Bosch durchlief ich die Abteilungen SIEM, Hadoop, Infrastruktur und Driving Experience. Dabei lernte ich, mich zügig in neue Themen und Teams einzuarbeiten. Zudem präsentierte ich regelmäßig technische Themen anforderungsgerecht für fachfremde Stakeholder. 

Im Masterstudium Machine Learning \& Data Analytics an der Hochschule Aalen vertiefte ich meine Kenntnisse in Deep Learning, Natural Language Processing und Computer Vision. Parallel arbeitete ich als Werkstudentin bei ZEISS SMT an Optimierungsverfahren für das Job-Shop-Scheduling-Problem sowie an der Analyse und dem Monitoring der Maschinenverfügbarkeit. Dabei verband ich analytische Methoden mit konkreten Fragestellungen der Praxis und stimmte mich eng mit den Fachabteilungen ab.

Bei Selfnet e.\,V., einer studentischen Initiative zum Betrieb der Internetinfrastruktur in Studentenwohnheimen, wirke ich an der Entwicklung des Mitgliederverwaltungssystems mit und sammle praktische DevOps-Erfahrung. Als Mentorin begleite ich neue Mitglieder beim Einstieg, erkläre technische Zusammenhänge verständlich und unterstütze bei Netzwerkproblemen. Diese Eigeninitiative und Teamarbeit möchte ich in die LBBW einbringen.

Besonders interessiert mich die Richtung Software-Development \& Innovation, ich bin aber offen dafür, auch andere Bereiche der IT kennenzulernen. Ich möchte lernen, wie IT-Lösungen für Kundinnen und Kunden einer Bank entstehen, und mir im Programm gezielt eine passende Zielposition erarbeiten.

Ich freue mich darauf, Sie in einem persönlichen Gespräch von meiner Motivation zu überzeugen und mehr über die Stationen des Programms zu erfahren.

Mit freundlichen Grüßen

\vspace{6mm}
Natalie Böhm

\end{document}